\section{Experiments}
\label{sec:experiments}

We evaluate MOMO at the workflow level: not only end-to-end accuracy but tool waste, context pollution, source-switch churn, partial-success handling, and cost. The control-plane claim is that these workflow-level pathologies, not retrieval, dominate single-agent tabular QA failures, and that contract-bounded executors plus runtime guards reduce them in measurable, mechanism-specific ways. Numeric cells dependent on result CSVs that have not yet been produced are marked \texttt{--}; the failure-mode taxonomy (Table~\ref{tab:failure}) is the qualitative contribution of this draft and is independent of pending numbers.

\subsection{Research Questions}

\begin{itemize}
\item \textbf{RQ1.} Does the control plane reduce \emph{workflow-level pathologies} --- turn waste, context-token growth, source-switch churn --- over a matched single-agent baseline under matched budgets and tool surfaces?
\item \textbf{RQ2.} Which failure modes from Table~\ref{tab:failure} does each control-plane mechanism (\texttt{results}, \texttt{cot}, \texttt{sprint}, \texttt{delegation}) reduce, and which remain?
\item \textbf{RQ3.} How does the subagent budget trade off cost, partial-rate, and end-to-end success, i.e., the cost-aware-decision axis?
\end{itemize}

\subsection{Setup}

Benchmark: the LakeQA-derived multi-source public-data task suite~\cite{wang2026lakeqa}, identical to the companion SANA paper's task set. Models: \texttt{gpt-5-mini} and \texttt{gpt-5.4-nano}. Tool budgets and tracing are shared across all conditions. The harness records, per task and per condition, the workflow-level evaluation surface: per-tool-call latency, token usage, cost, source-switch counts, peek-without-query rate, query-failure rate, context-token growth between turns, subagent budget-exhaustion rate, and per-contract status distribution. Search modes (\texttt{standard}, \texttt{ideal}, \texttt{preloaded}) are inherited from the SANA harness so MOMO conditions can be paired with each diagnostic search regime.

\subsection{Main Comparison}

Table~\ref{tab:main} compares single-agent baselines against MOMO variants on both accuracy-style metrics and the workflow-level differentiators (average source switches, average context tokens at submit). The control-plane prediction is that MOMO variants will not necessarily dominate on accuracy but will reduce source switches and context tokens at submit while maintaining or improving success rate at comparable cost.

\begin{table*}[t]
\caption{Main comparison across variants. Numeric cells pending result CSV. ``$\ast$'' indicates SANA-paper diagnostic conditions reused as references. Source-switch and context-token columns are the workflow-level differentiators.}
\label{tab:main}
\centering
\begin{tabular}{lllrrrrrrr}
\toprule
Variant & Search & Plan & Acc. & Success & Avg.\ calls & Avg.\ cost & Timeout & Src.\ sw. & Ctx.\ tok. \\
\midrule
Baseline single agent$^{\ast}$           & standard  & standard & -- & -- & -- & -- & -- & -- & -- \\
Ideal-search single agent$^{\ast}$       & ideal     & standard & -- & -- & -- & -- & -- & -- & -- \\
Preloaded single agent$^{\ast}$          & preloaded & standard & -- & -- & -- & -- & -- & -- & -- \\
MOMO-Results                             & preloaded & standard & -- & -- & -- & -- & -- & -- & -- \\
MOMO-Sprint (commitment)                 & preloaded & standard & -- & -- & -- & -- & -- & -- & -- \\
MOMO-Delegate                            & preloaded & standard & -- & -- & -- & -- & -- & -- & -- \\
MOMO-Delegate + Results                  & preloaded & standard & -- & -- & -- & -- & -- & -- & -- \\
\bottomrule
\end{tabular}
\end{table*}

\subsection{Search-Regime $\times$ Control-Plane Diagnostic}

Table~\ref{tab:stopping} pairs each search regime with each MOMO control-plane strategy. Standard search tests whether the agent can decide \emph{when} to stop searching. Preloaded sources removes source discovery entirely, exposing pure execution-discipline behavior. MOMO delegation bounds search and inspection separately, so the planner cannot burn the entire run on source thrash even under standard search.

\begin{table}[t]
\caption{Search-regime $\times$ control-plane diagnostic. Cells pending result CSV.}
\label{tab:stopping}
\centering
\begin{tabular}{lrrr}
\toprule
Strategy & Standard & Preloaded & Ideal \\
\midrule
None (baseline)           & -- & -- & -- \\
Sprint (commitment)       & -- & -- & -- \\
Delegation (MOMO)         & -- & -- & -- \\
\bottomrule
\end{tabular}
\end{table}

\subsection{Failure-Mode $\rightarrow$ Control-Plane Mechanism}

Table~\ref{tab:failure} is the qualitative spine of this paper. Each row enumerates a failure mode from \S\ref{sec:architecture}'s motivation, the trace signal we use to label it, the control-plane mechanism that addresses it, and the workflow metric in Table~\ref{tab:main} the row predicts to move. The mapping is independent of pending numeric results.

\begin{table*}[t]
\caption{Failure-mode taxonomy, addressing control-plane mechanism, and the workflow metric predicted to move under the addressing mechanism.}
\label{tab:failure}
\centering
\begin{tabular}{p{0.16\textwidth}p{0.30\textwidth}p{0.28\textwidth}p{0.18\textwidth}}
\toprule
Failure mode & Trace signal & Control-plane mechanism & Workflow metric impacted \\
\midrule
Over-searching          & repeated search calls after relevant source already found & bounded \texttt{search\_subagent} budget                          & avg.\ tool calls, cost \\
Schema thrashing        & repeated peek/query cycles on same file with no progress  & bounded \texttt{inspect\_subagent} budget; \texttt{cot} post-tool & avg.\ tool calls, success \\
File-handle errors      & dataset ID without exact path; ``file not found'' tracebacks & \texttt{\_FileReferenceGuard}                                  & timeout/ctx, success \\
Source drift            & inspect-time peeks/queries outside intended family        & \texttt{\_InspectSourceGuard}                                     & source switches \\
Zero-row loop           & repeated filtered queries returning zero rows             & worker-local retry; compact \texttt{failed} return                & avg.\ tool calls, cost \\
Context bloat           & schemas / row dumps / tracebacks dominate planner context & context firewall (compact returns)                                & ctx.\ tokens, success \\
Redundant verification  & re-queries of already-established facts                   & \texttt{sprint} \texttt{commitment\_reflection}; planner summaries & avg.\ tool calls \\
Premature submit        & final answer before \texttt{required\_outputs} covered    & \texttt{missing\_outputs} list in partial-status protocol         & success, accuracy \\
\bottomrule
\end{tabular}
\end{table*}

Per-row reduction --- the fraction of single-agent failures in each row, paired with the fraction under each MOMO configuration --- is the most direct measure of which mechanism contributes to which failure mode, and is the per-model breakdown the result CSVs will populate.

\subsection{Subagent-Budget Sensitivity}

Table~\ref{tab:budget} sweeps \texttt{max\_inspect\_subagent\_calls} with \texttt{max\_search\_subagent\_calls=3} and grace calls $= 1$ held fixed. The trade-off of interest is between tighter inspect budgets (lower cost, higher partial / budget-exhausted rates) and wider inspect budgets (lower partial rate but higher cost, with diminishing returns past a setting-specific threshold).

\begin{table}[t]
\caption{Subagent-budget sensitivity. Cells pending result CSV.}
\label{tab:budget}
\centering
\begin{tabular}{rrrrr}
\toprule
Inspect budget & Avg.\ subagent calls & Partial rate & Budget-exhausted rate & Cost \\
\midrule
4   & -- & -- & -- & -- \\
8   & -- & -- & -- & -- \\
12  & -- & -- & -- & -- \\
\bottomrule
\end{tabular}
\end{table}

\subsection*{Honesty Statement}

We report the harness, the variants, the workflow-level metrics, and the failure-mode $\rightarrow$ control-plane mechanism mapping. Numeric cells in Tables~\ref{tab:main}, \ref{tab:stopping}, and \ref{tab:budget} will be filled once the result CSVs are produced; the discussion that follows treats anticipated patterns as hypotheses, not claims.
